%! TeX root = ../thesis.tex
\chapter{Introduction}

Many body physics can give rise to interesting phenomena,
e.g., Bose-Einstein condensates (BEC) or superconductivity.
The general problem however
is the exponential increase of configurations with each particle added.
One approach to overcome this
is dynamical mean-field theory (DMFT)~\cite{Georges1996},
an approximation that is controlled in $1/z$ ($z$: number of neighbors)
and which successfully describes
Kondo physics~\cite{Wilson1975, Andrei1983}
and Mott insulators~\cite{MOTT1968}.
In DMFT, we can map a crystal lattice onto an
Anderson impurity model (AIM)~\cite{Anderson1961},
which we can solve among other things with
quantum Monte Carlo (QMC)~\cite{Hirsch1986, Gull2011}
or exact diagonalization (ED)~\cite{Caffarel1994} methods.
The disadvantages of the former compared to the latter are
that QMC is computationally demanding for low temperatures and the
analytical continuation from imaginary frequencies
is an ill-conditioned problem~\cite{Silver1990}.
The problem with ED is that it is only feasible for a few bath sites
which works well for isolators but not for materials with a Kondo peak.
Since we are interested in the ground state of a system and states close
to it,
we can limit the number of excitations,
using configuration interaction (CI)~\cite{Nesbet1967, Nesbet1970}.

The goal of this thesis is building a framework which can reproduce the work
of \citeauthor*{Lu2019}~\cite{Lu2019}.

\section{Anderson impurity model}
\label{sec:AIM}

Since no analytical solution of the Hubbard model~\cite{Hubbard1964} in
arbitrary dimensions exists,
we map it within DMFT onto a single-site AIM
to describe our system with fewer degrees of freedom.
In DMFT we hence ignore all spatial fluctuations
and only keep time correlations.
This approximation becomes exact in the limit of
infinite lattice coordination~\cite{Metzner1989}.
Our system contains a single impurity $i$ with energy level $\epsilon$
and on-site Coulomb interaction $U$.
This impurity is coupled by a hybridization $V_k$ to
bath sites with energies $\epsilon_k$
as illustrated in~\cref{subfig:geometry-AIM}.
Its Hamiltonian can be written as
\begin{equation}
    \label{eq:H_AIM}
    \begin{split}
        \ham_\text{AIM}
        = & \sum_\sigma \epsilon \fdag_\sigma \ff_\sigma
        + U \fdag_\up \fdag_\dn \ff_\dn \ff_\up           \\
          & + \sum_{k, \sigma} \epsilon_k \cdag\ks \cc\ks
        + \sum_{k, \sigma} \left(V_k \fdag_{\vphantom{k}\sigma} \cc\ks
        + V_k^* \cdag\ks \ff_{\vphantom{k}\sigma} \right),
    \end{split}
\end{equation}
with spin $\sigma \in \{\up, \dn\}$,
impurity creation (annihilation) operators $\fdag$ ($\ff$),
and bath creation (annihilation) operators $\cdag$ ($\cc$).
The first two terms of~\cref{eq:H_AIM} can be identified as the
impurity Hamiltonian $\ham_i$.

According to \citeauthor*{Lu2014}~\cite{Lu2014},
we can use different bath geometries
related to each other by unitary transformations,
as we want to represent the ground state of the system
with the fewest Slater determinants (SDs) possible.
Since it is not convenient to mix impurity and mean-field approximate sites,
only basis rotations within the mean-field sites are allowed,
preserving $\ham_i$.

The advantage of the representation in~\cref{subfig:geometry-AIM} is that bath
sites with low (high) onsite energy are
basically full (empty), indicated by their fill level.
Hopping is annotated by connecting sites by lines.
The disadvantage is that each bath site can interact with the impurity
and therefore is important.

We can make a basis transformation in which the impurity interacts with only
one bath site,
which in turn interacts with the next bath site etc., forming a chain.
This is shown in~\cref{subfig:geometry-tri}
where the single particle Hamiltonian can be expressed in a tridiagonal form.
The problem in this representation is that each bath site is partially occupied,
and the ground state requires an exponentially growing number of SDs.

The minimal number of SDs
to describe the ground state can be achieved in the natural orbital basis.
In this representation we create full (valence bath sites $v_i$)
and empty chains (conduction bath sites $c_i$)
which are coupled to the impurity and a mirror bath site $b$
as depicted in~\cref{subfig:geometry-nat}.
In order to be able to have an arbitrary filling $n$ of the impurity,
the mirror bath site has the opposite filling $1 - n$~\cite{Lu2014}.

\begin{figure}
    \centering
    \begin{subfigure}{0.3\textwidth}
        \centering
        %! TeX root = ../thesis.tex
\begin{tikzpicture}[
        inner sep=1mm,
        node distance=5mm,
    ]
    % draw sites
    \node [bath={0.8}, label=right:$b_1$] (b1) {};
    \node [bath={0.7}, label=right:$b_2$] (b2) [above=of b1] {};
    \node [bath={0.6}, label=right:$b_3$] (b3) [above=of b2] {};
    \node [bath half, label=right:$b_4$] (b4) [above=of b3] {};
    \node [bath={0.4}, label=right:$b_5$] (b5) [above=of b4] {};
    \node [bath={0.3}, label=right:$b_6$] (b6) [above=of b5] {};
    \node [bath={0.2}, label=right:$b_7$] (b7) [above=of b6] {};
    \node [impurity, label=left:$i$] (i) [left=of b4] {};
    % connect sites
    \foreach \i in {1,2,...,7}
    \draw (i.east) -- (b\i.west);
\end{tikzpicture}

        \caption{}
        \label{subfig:geometry-AIM}
    \end{subfigure}
    ~
    \begin{subfigure}{0.3\textwidth}
        \centering
        %! TeX root = ../thesis.tex
\begin{tikzpicture}[
        inner sep=1mm,
        node distance=5mm,
    ]
    % draw sites
    \node [impurity, label=right:$i$] (i) {};
    \node [bath half, label=right:$b_1$] (b1) [below=of i] {};
    \node [bath half, label=right:$b_2$] (b2) [below=of b1] {};
    \node [bath half, label=right:$b_3$] (b3) [below=of b2] {};
    \node [bath half, label=right:$b_4$] (b4) [below=of b3] {};
    \node [bath half, label=right:$b_5$] (b5) [below=of b4] {};
    % connect sites
    \draw (i.south) -- (b1.north);
    \foreach \i [remember=\i as \lasti (initially 1)] in {2,3,...,5}
    \draw (b\lasti.south) -- (b\i.north);
\end{tikzpicture}

        \caption{}
        \label{subfig:geometry-tri}
    \end{subfigure}
    ~
    \begin{subfigure}{0.3\textwidth}
        \centering
        \begin{tikzpicture}[
        inner sep=1mm,
        node distance=5mm,
    ]
    % draw sites
    \node [impurity, label=left:$i$] (i) {};
    \node [bath half, label=right:$b$] (b) [right=of i] {};
    \node [bath full, label=below:$v_1$] (v1) [below=of b] {};
    \node [bath empty, label=above:$c_1$] (c1) [above=of b] {};
    \foreach \i [remember=\i as \lasti (initially 1)] in {2,3,...,4}
    \node [bath full, label=below:$v_\i$] (v\i) [right=of v\lasti] {};
    \foreach \i [remember=\i as \lasti (initially 1)] in {2,3,...,4}
    \node [bath empty, label=above:$c_\i$] (c\i) [right=of c\lasti] {};
    % connect sites
    \draw (i.east) -- (b.west);
    \draw (i.south) -- (v1.north west);
    \draw (i.north) -- (c1.south west);
    \draw (b.south) -- (v1.north);
    \draw (b.north) -- (c1.south);
    \foreach \i [remember=\i as \lasti (initially 1)] in {2,3,...,4}
        {
            \draw (v\lasti.east) -- (v\i.west);
            \draw (c\lasti.east) -- (c\i.west);
        }
\end{tikzpicture}

        \caption{}
        \label{subfig:geometry-nat}
    \end{subfigure}
    \caption{
        Different bath geometries.
        Impurity $i$ interacting with bath sites $b_i$, or $b$, $v_i$, $c_i$.
        Site occupation is indicated by fill level.
        (adapted from~\cite{Lu2014}, Fig.~1)
    }
\end{figure}

We can further split the setup into two components as
shown in~\cref{fig:wave-function-split}.
The first component only includes the impurity, mirror bath site,
and $L$ valence and conduction bath sites,
with $L$ as a tunable parameter~\cite{Lu2019}.
In general this value can differ for
conduction/valence bath sites.
The second component includes all other bath sites,
which we want to approximate with CI.
This splits the Hamiltonian into three components
(see~\cref{fig:wave-function-split}):
$\ham_s$ which acts on the impurity and sites close to it,
$\ham_v$ which acts on the truncated chains,
and $\ham_\text{mix}$ which connects both.
For no excitation inside the truncated chains,
all valence bath sites $v_i$ are filled
and all conduction bath sites $c_i$ are empty which we write as
$\ket{\mathbb{1}_v} \otimes \ket{\mathbb{0}_c}$~\cite{Lu2019}.
In this work we will limit the number of excitations to singles and doubles,
meaning that up to one (two) occupations of spin-sites
can differ from the aforementioned state.

\begin{figure}[ht]
    \centering
    %! TeX root = ../thesis.tex
\begin{tikzpicture}[
        inner sep=1mm,
        node distance=5mm,
    ]
    % bit component
    \node [impurity, label=left:$i$] (i) {};
    \node [bath half, label=right:$b$] (b) [right=of i] {};
    \node [bath full, label=below:$v_1$] (v1) [below=of b] {};
    \node [bath full, label=below:$v_L$] (vL) [right=of v1] {};
    \node [bath empty, label=above:$c_1$] (c1) [above=of b] {};
    \node [bath empty, label=above:$c_L$] (cL) [right=of c1] {};
    \coordinate[node distance=20mm, right=of vL] (vhidden);
    \coordinate[node distance=20mm, right=of cL] (chidden);
    \node [node distance=14mm] (Hs) [left=of c1] {$\ham_s$};
    % vector component
    \node [bath full, label=below:$v_{L+1}$] (vL1) [right=of vhidden] {};
    \foreach \i [remember=\i as \lasti (initially 1)] in {2,3}
    \node [bath full] (vL\i) [right=of vL\lasti] {};
    \node [bath empty, label=above:$c_{L+1}$] (cL1) [right=of chidden] {};
    \foreach \i [remember=\i as \lasti (initially 1)] in {2,3}
    \node [bath empty] (cL\i) [right=of cL\lasti] {};
    \node [node distance=6mm] (Hv) [right=of cL3] {$\ham_v$};
    % % connect sites
    \draw (i.east) -- (b.west);
    \draw (i.south) -- (v1.north west);
    \draw (b.south) -- (v1.north);
    \draw(v1.east) -- (vL.west);
    \draw(c1.east) -- (cL.west);
    \draw (i.north) -- (c1.south west);
    \draw (b.north) -- (c1.south);
    \draw (vL.east) -- (vL1.west);
    \draw (cL.east) -- (cL1.west);
    \foreach \i [remember=\i as \lasti (initially 1)] in {2,3}
        {
            \draw (vL\lasti.east) -- (vL\i.west);
            \draw (cL\lasti.east) -- (cL\i.west);
        }
    \node (Hmix) [node distance=15mm, right=of b] {$\ham_\text{mix}$};
    % split box
    \node [
        inner sep=6mm,
        draw,
        dashed,
        rectangle,
        fit=(i) (b) (vL) (cL)
    ] {};
    \node [
        inner sep=6mm,
        draw,
        dashed,
        rectangle,
        fit=(vL1) (vL3) (cL1) (cL3)
    ] {};
\end{tikzpicture}

    \caption{
        Separating the Anderson Hamiltonian $\ham_\text{AIM}$
        into three components:
        $\ham_s$, $\ham_v$ act on sites enclosed by its corresponding box,
        $\ham_\text{mix}$ connects the two parts.
        (adapted from~\cite{Lu2019}, Fig.~2)
    }
    \label{fig:wave-function-split}
\end{figure}

\section{Fermions package}

We implement the concepts and ideas described in the previous section
in the module \textsc{Wavefunctions} of the
\textsc{Fermions}~\cite{Wallerberger2022} package,
which is written in Julia~\cite{Bezanson2017}.
First, we will give a short introduction of the package
and status quo,
which uses the occupation number basis for SDs.
Each site $k$ of the AIM contains two spin-sites $\sigma$
with occupation $n\ks \in \{0, 1\}$ for fermions.
Each SD $\ket{n_n, \ldots, n_3, n_2, n_1}$
can therefore be interpreted as a bit pattern of zeros and ones
represented by an unsigned integer
with creation/annihilation operators acting as bit flips.
The vacuum state $\ket{0}$ can be depicted as zero consisting of no particles.
To store the wave function,
we can use a dictionary with keys acting as SDs
and associated values acting as probability amplitudes
\begin{equation}
    \ket{\psi} = \sum_i \alpha_i \ket{S_i},
\end{equation}
which is currently implemented in the struct \mintil{Wavefunction}.
For normalization $\sum_i |\alpha_i|^2=1$ must be fulfilled.
To satisfy the fermionic anticommutator relations,
the annihilation operator is defined as
\begin{equation}
    \label{eq:anticommutator-definition}
    \begin{split}
        \cc_k \ket{\ldots, n_k=1, \ldots}
         & \coloneq \prod_{i=1}^{k-1} (-1)^{n_i} \ket{\ldots, n_k=0, \ldots} \\
         & = (-1)^{\sum_{i=1}^{k-1} n_i} \ket{\ldots, n_k=0, \ldots}.
    \end{split}
\end{equation}
